%%%%%
%
%% Dans ce fichier, les �ventuels accents sont cod�s
%% en ISO-8859-15 (aka latin9)
%% La ligne suivante donne trois exemples 
%% e accent aigu : �, e accent grave : �, e-dans-l'o : �
\section{Rotation des 3 composantes du vent autour des axes X et Y}

\itemstyle{Op\'erateur~:} \operfmt{Rotat\_XY}
\index{Rotat\_XY}

%%%%%OBJET%%%%
\itemstyle{Objet~:}
Effectue une double rotation des 3 composantes du vent autour des axes X et Y du rep�re de l'an�mom�tre.
Cette rotation est mise en oeuvre par exemple dans le cas de l'application du calcul d'un "Planar-Fit" voir 
\cite{Wilczak-2001}

%%%%%ARGUMENTS%%%%
\bigskip
%% Dans ce fichier, les �ventuels accents sont cod�s
%% en ISO-8859-15 (aka latin9)
%% La ligne suivante donne trois exemples 
%% e accent aigu : �, e accent grave : �, e-dans-l'o : �
%%%%%ENTREES%%%%
\begin{tabarg}
\multicolumn{3}{l}{\titretab {Entr\'ees}}\\
\hline
$U$ & Vecteur, \texttt{float} & Composante du vent suivant axe X \\
$V$ & Vecteur, \texttt{float} & Composante du vent suivant axe Y \\
$W$ & Vecteur, \texttt{float} & Composante du vent suivant axe Z \\
$\alpha$ & Coeff., \texttt{float} & Angle de rotation autour de l'axe X (degr�s)\\
$\beta$ & Coeff., \texttt{float} & Angle de rotation autour de l'axe Y (degr�s) \\
$ b_0$ & Coeff., \texttt{float} & Correction d'offset de $W$\\
\end{tabarg}

\medskip
%%%%%SORTIES%%%%
\begin{tabarg}
\multicolumn{3}{l}{\titretab {Sorties}}\\
\hline
$U_r$ & Vecteur, \texttt{float} & Composante U dans le nouveau rep�re (m\^eme unit� que $U$)\\
$V_r$ & Vecteur, \texttt{float} & Composante V dans le nouveau rep�re (m\^eme unit� que $V$)\\
$W_r$ & Vecteur, \texttt{float} & Composante W dans le nouveau rep�re (m\^eme unit� que $W$)\\
\end{tabarg}


%%%%%FORMULE%%%%
\itemstyle{Formule}

%% Formule de l'operateur : Rotat_XY
%% Dans ce fichier, les �ventuels accents sont cod�s
%% en ISO-8859-15 (aka latin9)
%% La ligne suivante donne trois exemples 
%% e accent aigu : �, e accent grave : �, e-dans-l'o : �

$$
\begin{cases}
U_r = \cos\alpha.U - \sin \alpha .\sin\beta.V + \sin\alpha.\cos \beta.(W - b_0) \\
V_r = \cos \beta .V           + \sin \beta.(W - b_0) \\
W_r = \sin \alpha.U - \cos \alpha .\sin \beta .V + \cos \alpha . \cos \beta .(W - b_0) \\
\end{cases}
$$



%%%%%AUTEUR%%%%
\itemstyle{Auteur~:}
CNRM/GMEI/TRAMM

%%%%%REFERENCE%%%%
\itemstyle{R\'ef\'erences~:}
Wilczak 2001 \cite{Wilczak-2001}
